\documentclass[12pt]{article}

% ── Packages ──────────────────────────────────────────────────────────────────
\usepackage[a4paper, margin=1in]{geometry}
\usepackage{setspace}
\usepackage{newtxtext, newtxmath}
\usepackage{titlesec}
\usepackage{fancyhdr}
\usepackage{hyperref}
\usepackage{xurl}
\usepackage{doi}
\usepackage{parskip}
\usepackage{tabularx}
\usepackage[natbibapa]{apacite}

% ── Spacing and indentation ───────────────────────────────────────────────────
\doublespacing
\setlength{\parindent}{0.5in}
\setlength{\parskip}{0pt}

% ── Section headers (APA 7th Level 1 and Level 2) ────────────────────────────
\titleformat{\section}
  {\normalfont\fontsize{12}{14}\bfseries\centering}{}{0em}{}
\titleformat{\subsection}
  {\normalfont\fontsize{12}{14}\bfseries}{}{0em}{}
\titlespacing*{\section}{0pt}{12pt}{0pt}
\titlespacing*{\subsection}{0pt}{12pt}{0pt}

% ── Header: page numbers top right ───────────────────────────────────────────
\pagestyle{fancy}
\fancyhf{}
\rhead{\thepage}
\renewcommand{\headrulewidth}{0pt}

\hypersetup{colorlinks=true, urlcolor=black, linkcolor=black}

% ──────────────────────────────────────────────────────────────────────────────
\begin{document}

% ── Title Page (no page number, not counted) ─────────────────────────────────
\begin{titlepage}
  \thispagestyle{empty}
  \begin{center}
    \vspace*{2in}
    \textbf{A Differentiable E-Z Reader: Integrating Neural Language Models\\
    into a Cognitive Model of Eye-Movement Control in Reading}

    \vspace{24pt}
    Member 1: Efe Kahraman (u384661)\\[6pt]
    Member 2: Lorenzo Cipriani (u343767) \\[6pt]
    Member 3: Arturo Dalbagno (u381361) \\[6pt]
    Member 4: Meda Kuizinaite (u702127)

    \vspace{24pt}
    {Research Workshop: CSAI}\\[4pt]
    {Nicenboim, Bruno}\\[4pt]
    {Date: 10/04/2026}
  \end{center}
\end{titlepage}

% ── Technology Statement Page ─────────────────────────────────────────────────
\newpage
\thispagestyle{empty}
\begin{center}
  {\large\textbf{Technology Statement}}
\end{center}
\vspace{12pt}

\begin{tabularx}{\textwidth}{|l|X|}
\hline
\textbf{1.} & \textbf{Did you use any tools or services to help in the prewriting
              stages of the writing process?} \\
\hline
            & Claude Code, Gemini CLI \\
\hline
\textbf{2.} & \textbf{Did you use any tools/services to paraphrase or rephrase
              your text (e.g., the Academic Phrasebank, Quillbot, DeepL, ChatGPT)?} \\
\hline
            & Cursor (Claude Sonnet 4.6) \\
\hline
\textbf{3.} & \textbf{Did you use any tools or services to check spelling or grammar?} \\
\hline
            & Cursor (Claude Sonnet 4.6) \\
\hline
\textbf{4.} & \textbf{Did you use any tools or services to put (parts of) the text
              into the required format (e.g., APA), including references
              (e.g., EndNote, Zotero, Scribbr)?} \\
\hline
            & Cursor (Claude Sonnet 4.6), Overleaf \\
\hline
\textbf{5.} & \textbf{Did you use any tools or services to generate part of the text?} \\
\hline
            & Cursor (Claude Sonnet 4.6) \\
\hline
\textbf{6.} & \textbf{Did you use any generative AI tools or software for other
              aspects of the assignment?} \\
\hline
            & Overleaf (for document formatting) \\
\hline
\end{tabularx}

\vspace{24pt}

\noindent\textbf{Confirmation Statement:} By submitting this assignment, we,
Efe Kahraman, Lorenzo Cipriani, Arturo Dalbagno, Meda Kuizinaite, confirm that:

\begin{enumerate}
  \item The assignment is our own intellectual property and that ideas as well
        as language from other sources have been properly cited, referenced,
        and quoted (when necessary). All language and ideas not our own have
        been properly identified as such.
  \item We have disclosed all technology that we used in the writing process.
  \item We claim responsibility for our submitted work.
\end{enumerate}

\newpage

% ── Body starts at page 1 ────────────────────────────────────────────────────
\setcounter{page}{1}
\pagestyle{fancy}

\begin{center}
  \textbf{A Differentiable E-Z Reader: Integrating Neural Language Models\\
  into a Cognitive Model of Eye-Movement Control in Reading}
\end{center}

\vspace{6pt}
\noindent\textbf{Group Members:}\;
Member 1: Efe Kahraman;
Member 2: Lorenzo Cipriani;
Member 3: Arturo Dalbagno;
Member 4: Meda Kuizinaite{\hspace{2.2cm}}

\vspace{12pt}

% ── Research Motivation ───────────────────────────────────────────────────────
\section{Research Motivation}

Eye movements during reading consist of rapid jumps between fixation points
(saccades) alternating with brief periods of relative stillness (fixations),
and these patterns are directly sensitive to the linguistic properties of the
text being processed. \citet{rayner1998eye} synthesized over two decades of
evidence showing that eye movement data reflect moment-to-moment cognitive
processing during reading, providing the empirical basis for word-level
measures such as first fixation duration (FFD), gaze duration, total reading
time (TRT), and skip rate as standard indices of lexical processing
difficulty.

The E-Z Reader model \citep{reichle1998toward} provides a computational
account of eye movement control during reading. The model proposes two
sequential processing stages for each word: a familiarity check (L1) that
initiates saccade programming toward the next word, and a full lexical access
stage (L2) that shifts attentional resources. Both stages are computed through
assumed parametric functions of word frequency and contextual predictability,
inputs that must be obtained from precomputed lexical databases and norming
studies.

A separate line of research has approached reading time prediction from an
information-theoretic perspective. \citet{levy2008expectation} formalized the
proposal that processing difficulty during reading is proportional to the
surprisal of a word, defined as its negative log-probability given the
preceding context. \citet{goodkind2018predictive} demonstrated empirically
that the predictive power of surprisal for reading times scales as a linear
function of language model quality, directly linking improvements in language
modeling to gains in behavioral prediction. Although these approaches achieve
strong predictive performance at the word level, they reduce the model's output
to a single scalar per word, which does not permit decomposition of processing
difficulty into distinct cognitive stages.

The present study addresses a gap that neither tradition has closed. The
cognitive parameters of E-Z Reader are typically calibrated by grid search
against group-averaged fixation distributions, because the simulation is
non-differentiable and admits no gradient-based optimization. This restricts
parameter recovery to coarse aggregate behavior and prevents joint training
with modern context-sensitive predictability estimates obtained from neural
language models. \citet{lopesrego2024language} recently replaced cloze norms
with LLM-derived predictability inside the OB1-reader cognitive model,
demonstrating empirical gains, but their pipeline preserved the
non-differentiable simulation, so the cognitive parameters could not be
co-optimized with the language model encoder. We propose a differentiable
reformulation of E-Z Reader that admits gradient-based fitting of cognitive
parameters at the word level, jointly with a neural language model encoder
trained end-to-end on human eye-tracking data.

\noindent\textbf{Scientific relevance.} This work tests whether E-Z Reader's
stage decomposition (familiarity check versus lexical access) survives
end-to-end optimization with a neural encoder. If the gradient-fit cognitive
parameters remain interpretable and comparable in magnitude to those reported
by \citet{reichle1998toward}, the result provides quantitative evidence that
the dual-stage architecture captures real psychological structure rather than
reflecting hand-tuning artifacts of grid-search calibration. If the parameters
drift toward a better-fitting but theoretically distinct configuration, the
divergence informs future cognitive theories of eye-movement control. More
broadly, eye movements during reading are a non-invasive behavioral window
onto cognitive processing, and constraining the computational architecture
that produces them informs models of how the brain coordinates lexical
recognition with visual-motor control during reading.

\noindent\textbf{Societal relevance.} Word-level reading-time prediction
underpins applications in reading-difficulty estimation, automatic text
simplification, accessibility assessment for second-language learners, and
atypical-reading research more broadly. Models that combine quantitative
accuracy with interpretable cognitive stages are more defensible in clinical
and educational deployments than black-box regression models, because
stage-level decomposition can localize where in the reading process ---
familiarity check or lexical access --- a particular text or reader
population is encountering difficulty.

% ── Research Question ─────────────────────────────────────────────────────────
\section{Research Question}

Can a differentiable reformulation of E-Z Reader, with cognitive parameters
and a neural language model encoder fit jointly by gradient descent on
word-level eye-tracking data, achieve reading-time prediction accuracy
comparable to direct neural regression baselines that predict reading times
from raw text without any cognitive architecture --- specifically (i) a
matched-capacity direct-regression model that uses the same Llama-3.2-1B
encoder with four independent output heads predicting FFD, gaze duration,
TRT, and skip rate, and (ii) the leading systems from the CMCL 2021 Shared
Task on Eye-Tracking Prediction \citep{hollenstein2021cmcl}, including the
fine-tuned RoBERTa-large model of \citet{liohio2021cmcl} and the RoBERTa-base
model with task-adaptive pretraining of \citet{litorontocl2021cmcl} ---
while preserving the mechanistic interpretability of the underlying cognitive
architecture?

% ── Background ────────────────────────────────────────────────────────────────
\section{Background}

The E-Z Reader model \citep{reichle1998toward} accounts for eye movement
patterns across a range of reading conditions by decomposing word processing
into two sequential cognitive events. The familiarity check (L1) is completed
rapidly and triggers saccade programming toward the next word. Lexical access
(L2) is slower and governs the shift of focal attention. The model generates
predictions at the level of aggregated data, fitting mean fixation durations
and fixation probabilities across groups of words or experimental conditions.
Because its inputs must be supplied from externally normed databases, the model
cannot be applied to arbitrary text, and its parameters are typically
calibrated by hand to fit a specific corpus.

The theory of surprisal \citep{levy2008expectation} provides a complementary
framework for understanding lexical processing difficulty. Under this account,
the processing cost of a word is proportional to its negative log-probability
given the preceding context, connecting cognitive processing difficulty directly
to probabilistic language modeling. \citet{goodkind2018predictive} formalized
this connection empirically, showing that the predictive power of surprisal for
reading times is a linear function of language model quality.
\citet{wilcox2023testing} extended this finding across eleven languages,
confirming that the relationship between language model performance and
behavioral reading time prediction generalizes beyond English.

Despite this predictive success, surprisal-based approaches do not model the
cognitive processes that mediate the relationship between word probability and
fixation behavior, yielding a single scalar prediction per word without
decomposing processing difficulty into sequential stages. A recent attempt to
bridge this divide replaced cloze predictability norms with LLM-derived
probabilities as inputs to OB1-reader, a cognitive model of eye-movement
control, and found that LLM-based predictability provides a better fit to
human eye-tracking data than traditional cloze norms
\citep{lopesrego2024language}. However, that approach preserves the
non-differentiable simulation architecture and does not allow the cognitive
model parameters to be jointly optimized with the language model encoder.

The present study addresses this limitation directly. E-Z Reader is implemented
as a discrete event simulation whose operations are not differentiable,
precluding gradient-based optimization and joint training with a learned
encoder. By reformulating E-Z Reader as a differentiable computation graph and
replacing its parametric input functions with the output of a neural language
model, both components can be trained simultaneously on human eye-tracking
data, enabling the model to learn cognitively structured predictions without
relying on external norms.

% ── Methods ───────────────────────────────────────────────────────────────────
\section{Methods}

\subsection{Human Subject Experiment}

The present study will make use of two existing eye-tracking corpora as the
source of behavioral data. The Ghent Eye-Tracking Corpus
\citep[GECO;][]{cop2017presenting} consists of recordings from 14 monolingual
English participants reading a complete novel, yielding approximately 54,000
word-level observations. The Provo Corpus \citep{luke2018provo} contains
eye-tracking data from 84 participants reading 55 short expository passages
(2,654 words in total), accompanied by cloze-based contextual predictability
norms collected from a separate norming sample. Both corpora employed a
within-subject design in which all participants read the same materials. The
dependent variables are first fixation duration (FFD), gaze duration, total
reading time (TRT), and skip rate, measured at the level of individual words.
The independent variables are word-level properties that naturally vary across
stimuli: orthographic word length, word frequency, and contextual
predictability. Behavioral measures will be aggregated across participants per
word position and analyzed using word-level correlation metrics, serving as
both training targets and evaluation criteria for the machine learning
experiment.

\subsection{Machine Learning Experiment}

The GECO corpus will be used for model training and in-distribution evaluation,
with sentences partitioned into training, validation, and test sets at the
level of complete texts to prevent data leakage. The Provo corpus will be held
out entirely as a cross-corpus generalization benchmark, allowing evaluation on
text of a different genre and from a different participant population than those
seen during training.

Four models will be implemented and compared. The first is the original E-Z
Reader \citep{reichle1998toward}, using externally supplied word frequency and
predictability values as inputs, serving as the cognitive baseline. The second
is a differentiable reformulation of E-Z Reader using the same formula-based
inputs but with all model parameters optimized via gradient descent, isolating
the contribution of end-to-end parameter learning. The third is the full
proposed model: a pretrained causal language model serves as an encoder,
producing contextual word representations that are projected through three
learned output heads predicting L1 processing time, L2 processing time, and
skip probability, which are then passed into the differentiable E-Z Reader to
produce FFD, gaze duration, TRT, and skip rate predictions. The fourth model is
a direct neural regression baseline using the same encoder but with four
independent output heads predicting each reading time measure directly, without
any cognitive architecture, serving as the critical control for assessing
whether the cognitive structure provides any benefit over unconstrained
regression.

Model performance will be evaluated using word-level Pearson correlation
between predicted and observed values for each dependent variable, computed
separately on the GECO held-out test set and on the full Provo corpus.
Cross-corpus generalization will be assessed by comparing performance across
the two evaluation sets for all four models.

% ── References ────────────────────────────────────────────────────────────────
\newpage
\bibliographystyle{apacite}
\bibliography{references}

\end{document}
